\section{A smaller sign-algebra trace and its exact source correspondence}
\label{sec:compressed-permanent-trace}

The permanent trace of Section~\ref{sec:permanent-trace-projection}
admits a second, smaller finite-algebra representation. We retain the
same $n^2$ coordinates $Z_{ij}$ and the same permanent polynomial;
we construct the correspondence between the two representations explicitly.
The sign sum used below is the Ryser--Glynn expression recorded in
\cite[Equation (8)]{LandsbergRessayre2017SignTrace}; the original
Glynn article is \cite{Glynn2010SignTrace}. All identities and maps
needed here are proved below, independently of those references.

\subsection{The sign algebra, its integral basis, and its regular trace}

Fix an integer $n\ge1$. Put $O=\mathbb Z[1/2]$, $R=O[Z_{ij}]$,
$r_n=2^{n-1}$, and
\[
 B_n=O[d_2,\ldots,d_n]/(d_j^2-1:2\le j\le n),\qquad d_1=1.
\]
For $n=1$ this means $B_1=O$. For each subset
$S\subseteq\{2,\ldots,n\}$ write $d_S=\prod_{j\in S}d_j$,
including $d_\varnothing=1$. The elements $d_S$ form an $O$-basis.
Indeed, successive division by the monic polynomials $d_j^2-1$
gives a unique remainder with exponents zero or one: viewing the
ring as a polynomial ring in the last variable over the preceding
quotient, monic division gives a free module with basis $1,d_j$.
Induction proves the assertion over $O$ and equally over $\mathbb Z$.
Multiplication is
\[
 d_Sd_T=d_{S\mathbin{\triangle}T},
\]
where $S\mathbin{\triangle}T=(S\setminus T)\cup(T\setminus S)$.
There are exactly $r_n$ basis elements, one for each independent
membership choice in a set of size $n-1$.

Write $\tau_B(b)=\operatorname{Tr}_{B_n/O}(m_b)$ for the regular
trace, where $m_b(x)=bx$. Multiplication by $d_S$ permutes the basis
by $T\mapsto S\mathbin{\triangle}T$. A basis vector is fixed if
and only if $S=\varnothing$, since cancelling $T$ under symmetric
difference gives that equivalence. Thus
\begin{equation}\label{eq:compressed-regular-trace}
 \tau_B\Bigl(\sum_S b_Sd_S\Bigr)=r_n b_\varnothing.
\end{equation}
This identity holds first over $\mathbb Z$ and then over every
coefficient ring by taking the same finite diagonal sum.

There is also an explicit split basis. Let
$\Delta_n=\{\delta\in\{-1,1\}^n:\delta_1=1\}$ and set
\[
 \epsilon_\delta=\prod_{j=2}^n\frac{1+\delta_jd_j}{2}.
\]
Each factor is idempotent, since
$(1+\delta_jd_j)^2/4=(1+\delta_jd_j)/2$.
Opposite factors multiply to $(1-d_j^2)/4=0$, and summing the
two factors for one coordinate gives one. Multiplying these identities
shows that the $\epsilon_\delta$ are orthogonal idempotents summing
to one. Evaluation at $d_j=\eta_j$ gives
$\epsilon_\delta(\eta)=1$ when $\eta=\delta$ and zero otherwise.
The identity
$d_S=\sum_\delta(\prod_{j\in S}\delta_j)\epsilon_\delta$
follows by multiplying $d_j\epsilon_\delta=\delta_j\epsilon_\delta$
and summing the idempotents. It proves the mutually inverse algebra maps
\[
 B_n\longrightarrow O^{\Delta_n},\quad b\longmapsto(b(\delta))_\delta,
 \qquad(c_\delta)_\delta\longmapsto\sum_\delta c_\delta\epsilon_\delta.
\]
In particular $\tau_B(b)=\sum_\delta b(\delta)$.

For completeness the factor inverted here has exact integral content.
Before localizing, the evaluation matrix from the $d_S$ basis to
the integral point basis is
$H_{\delta,S}=\prod_{j\in S}\delta_j$. For two subsets $S,T$,
the sum $\sum_\delta H_{\delta,S}H_{\delta,T}$ is $r_n$ if
$S=T$ and zero otherwise: in the latter case choose a coordinate
in $S\mathbin{\triangle}T$ and pair each $\delta$ with the sign
vector obtained by flipping that coordinate. Therefore
\[
 H^{\mathsf T}H=r_n I,\qquad H^{-1}=r_n^{-1}H^{\mathsf T}
 \quad\text{over }\mathbb Q.
\]
The image of $\mathbb Z[d_2,\ldots,d_n]/(d_j^2-1)$ in
$\mathbb Z^{\Delta_n}$ consists exactly of integer vectors $v$
such that every component of $H^{\mathsf T}v$ is divisible by $r_n$:
necessity follows from the first identity, and sufficiency follows
by applying the displayed inverse. Thus the integral evaluation map
is injective and becomes an isomorphism after inverting two. For
$n\ge2$ it is not surjective: the coordinate vector supported at
a single $\delta$ has $H^{\mathsf T}v$ with every component $1$
or $-1$, which is not divisible by $r_n\ge2$.
No integral identification of these two lattices has been made.

\subsection{The exact parity computation, with both orientations}

Put $\chi_n=\prod_{j=1}^n d_j=d_{\{2,\ldots,n\}}$ and define
row-linear elements and their multiplier by
\begin{equation}\label{eq:compressed-row-multiplier}
 L_i(Z)=\sum_{j=1}^n d_j Z_{ij},\qquad
 v_n(Z)=r_n^{-1}\chi_n\prod_{i=1}^n L_i(Z)
 \quad\text{in }B_n\otimes_O R.
\end{equation}

\begin{theorem}\label{thm:compressed-permanent-trace}
The regular multiplication trace on the rank-$r_n$ algebra satisfies
\[
 \operatorname{Tr}_{B_n\otimes R/R}(m_{v_n(Z)})
 =\sum_{\sigma\in\mathfrak S_n}\prod_{i=1}^n Z_{i,\sigma(i)}
 =\operatorname{per}_n(Z).
\]
The identity holds over $O$ and after any homomorphism from $O$
to a commutative ring, in particular over $\mathbb Q$ and every
field of odd characteristic.
\end{theorem}
\begin{proof}
Expand the product in \eqref{eq:compressed-row-multiplier} without
changing its row order:
\[
 v_n(Z)=r_n^{-1}
 \sum_{(j_1,\ldots,j_n)\in[n]^n}
      \chi_n d_{j_1}\cdots d_{j_n}\prod_{i=1}^n Z_{ij_i}.
\]
For a tuple let $m_j$ be the number of its entries equal to $j$.
The group-basis monomial $\chi_n\prod_i d_{j_i}$ is one exactly
when $m_j+1$ is even for each $j\ge2$, or equivalently when each
of $m_2,\ldots,m_n$ is odd. In that case write $m_j=1+2a_j$
for $j\ge2$, with $a_j\ge0$. The equality $\sum_jm_j=n$ becomes
\[
 m_1+2\sum_{j=2}^n a_j=1.
\]
It forces $m_1=1$ and every $a_j=0$. Conversely, if every column
is chosen once, all the exponents are even as required. Thus exactly
the permutation tuples contribute to the coefficient of $d_\varnothing$,
each with coefficient $r_n^{-1}$ times its original row monomial.
Equation~\eqref{eq:compressed-regular-trace} multiplies that coefficient
by the exact factor $r_n$, proving the formula. For $n=1$ the
same equation says $m_1=1$; directly $r_1=\chi_1=d_1=1$ and
$v_1(Z)=Z_{11}$ in the rank-one algebra. Finally all the calculations
are identities in the displayed free algebra over $O$, so applying
a coefficient homomorphism preserves them, including their traces.
\end{proof}

The split-basis form of this equality is the fully signed expression
\[
 \operatorname{per}_n(Z)=2^{1-n}
 \sum_{\delta\in\Delta_n}
 \Bigl(\prod_{j=1}^n\delta_j\Bigr)
 \prod_{i=1}^n\Bigl(\sum_{j=1}^n\delta_j Z_{ij}\Bigr).
\]
The prefactor and product of signs are retained. To give the other
orientation explicitly, use a second copy of the sign algebra with
generators $e_2,\ldots,e_n$, $e_1=1$, and put
\[
 \widehat v_n(Z)=2^{1-n}
 \Bigl(\prod_{i=1}^n e_i\Bigr)
 \prod_{j=1}^n\Bigl(\sum_{i=1}^n e_i Z_{ij}\Bigr).
\]
Let the coefficient-ring involution send $Z_{ij}$ to $Z_{ji}$,
and send $d_j$ in the first sign algebra to $e_j$ in the second.
This is an isomorphism of the underlying $O$-algebras with polynomial
coefficients; its square, with the reverse generator renaming, is
the identity. It sends $v_n(Z)$ to $\widehat v_n(Z)$ by direct
substitution. The trace of the latter is
$\operatorname{per}_n(Z^{\mathsf T})=\operatorname{per}_n(Z)$:
the term $\prod_i Z_{\sigma(i),i}$ becomes
$\prod_j Z_{j,\sigma^{-1}(j)}$, and inversion permutes
$\mathfrak S_n$ bijectively. Thus transposition is a specified
coordinate map, not an implicit interchange of the original rows
and columns in \eqref{eq:compressed-row-multiplier}.

Over $\mathbb Z$ the unscaled numerator
$w_n(Z)=\chi_n\prod_iL_i(Z)$ is an element of the integral group
algebra and satisfies $\operatorname{Tr}(m_{w_n})=r_n\operatorname{per}_n$.
For $n\ge2$ the scaled element $v_n$ is not in that integral
group algebra with polynomial coefficients. The monomial
$Z_{11}\cdots Z_{n1}$ occurs exactly once in the expansion and has
coefficient $\chi_n/r_n$, which is not integral in the $d_S$ basis.
The resulting trace polynomial nevertheless belongs to $\mathbb Z[Z]$
by the proved identity. It can be reduced to characteristic two
through that integer polynomial or the original split integral trace
matrix. There is no homomorphism $O\to k$ for a field $k$ of
characteristic two, because the unit two would map to zero.
The scaled sign-algebra operator is therefore not specialized by
such a nonexistent coefficient map. Its integer numerator has trace
zero there for $n\ge2$, which cannot by itself recover the permanent;
at $Z=I_n$ the permanent is one.

\subsection{The original source fibre over the same coefficient ring}

Retain the original map with its source orientation $(x,y,w)$ and
target orientation $(a,b,c)$:
\begin{align*}
 F_1&=(1+xy)^3w+y^2(1+xy)(4+3xy),\\
 F_2&=y+3x(1+xy)^2w+3xy^2(4+3xy),\\
 F_3&=2x-3x^2y-x^3w.
\end{align*}
It still has $\det DF=-2$, and its retained polynomial and derivative
are $P=cr^3-2r^2+br-2a$ and $P'=3cr^2-4r+b$. Neither expression
is altered by the sign algebra. The original fibre at $(-1/4,0,0)$
is defined over $O$, and the isomorphism
\[
 O[x,y,w]/(F_1+1/4,F_2,F_3)
 \longrightarrow O[X]/(X^3-X),\qquad
 (x,y,w)\longmapsto\bigl(X,-3X/2,(27X^2-1)/4\bigr)
\]
continues to hold there. To check the base ring explicitly, put
$h=2-3xy-x^2w$. The identity
$1=h/2+x(3y+xw)/2$ makes $(x)$ and $(h)$ comaximal after imposing
$F_1+1/4,F_2$. As $F_3=xh$, the quotient splits into its $x=0$
and $h=0$ factors by the same two-ideal intersection argument.
The first is $O$, with $(x,y,w)=(0,0,-1/4)$. In the second,
$x(3y+xw)=2$ makes $x$ a unit; putting $r=y+1/x$, $\alpha=1/x$
gives $F_1=r^2+r\alpha$, $F_2=4r+2\alpha$, hence
$\alpha=-2r$, $r^2=1/4$. Conversely these equations make $r$
a unit and recover the original source coordinates and target
equations. This factor is $O[r]/(r^2-1/4)$, with
$x=-2r$, $y=-3x/2$, $w=13/2$. All inverses used are powers of
two or the stated source unit. The two factors give exactly
$O[X]/(X^3-X)$, with its original $x=0$ point retained.

The idempotents $e_-=(X^2-X)/2$, $e_0=1-X^2$,
$e_+=(X^2+X)/2$ give evaluation at $X=-1,0,1$ over $O$ as
in Proposition~\ref{prop:sv-full-fibre}. Take $N=3n^2$ tensor
factors and write $\mathcal A_N$ for the resulting full source
algebra over $O$. Its point set is $\Omega_N=\{-1,0,1\}^N$;
write $E_\omega$ for its product idempotents. It has rank $3^N$,
and $\Lambda_N=\bigoplus_{\omega\in\Omega_N}\mathbb Z E_\omega$
is the original split integral lattice, with its coordinatewise
multiplication. In particular $\mathcal A_N=\Lambda_N\otimes O$.
This lattice differs from the unsplit integral presentation with
source-coordinate denominators, as already proved for each factor.

For a permutation $\sigma$, let $\omega(\sigma)$ be the original
source point with $B_{ij}=1$ exactly when $j=\sigma(i)$ and with
the unique row and column prefix flags of the preceding construction.
In every witness coordinate $h$ this point has
\[
 X_h=2B_h-1,\qquad y_h=-3X_h/2,\qquad
 w_h=(27X_h^2-1)/4.
\]
Here $B_h$ denotes the appropriate original witness bit, including
the two groups of flags. Distinct permutations have distinct
$B$-arrays and therefore distinct points. Put
\[
 f_n=\sum_{\sigma\in\mathfrak S_n}E_{\omega(\sigma)},
 \qquad A_{\mathrm{perm}}=f_n\mathcal A_N.
\]
Orthogonality proves $f_n^2=f_n$ and gives the algebra decomposition
$\mathcal A_N=A_{\mathrm{perm}}\oplus(1-f_n)\mathcal A_N$.
The first ideal has identity $f_n$ and rank $n!$. The map
$a\mapsto f_na$ is a surjective algebra map to this ideal with
kernel $(1-f_n)\mathcal A_N$. It retains the complementary ideal
as its explicitly specified kernel, including every boundary point
with some $X_h=0$.

The projected original multiplier is exactly
\begin{equation}\label{eq:compressed-original-u}
 u_n(Z)=\sum_{\sigma\in\mathfrak S_n}
           E_{\omega(\sigma)}\prod_{i=1}^n Z_{i,\sigma(i)}.
\end{equation}
Indeed the unweighted constraint-clause product is one precisely on those uniquely
flagged assignments, each appended weighted clause supplies the
specified entry, and the original product idempotent kills every
boundary factor. These are precisely the evaluations of the
projected multiplier established above; the split evaluation
isomorphism proves equality as an element, not only as a trace.
Since $u_n=f_nu_n$, its multiplication is zero on the full
complementary ideal. Consequently
\[
 \operatorname{Tr}_{\mathcal A_N}(m_{u_n})
 =\operatorname{Tr}_{A_{\mathrm{perm}}}(m_{u_n})
 =\operatorname{per}_n(Z).
\]
The full original rank $3^{3n^2}$ and all its zero eigenvalues
remain present in the first trace.

\subsection{A common tensor module and an exact sequence of traces}

Let $T_n=(O^n)^{\otimes n}$, with factor $i$ assigned to row $i$.
For $\boldsymbol j=(j_1,\ldots,j_n)\in[n]^n$ denote its tensor
basis vector by $t_{\boldsymbol j}$. Define three $O$-linear maps
on every basis vector by
\[
 \begin{split}
 S_n:T_n&\longrightarrow A_{\mathrm{perm}},\qquad
 S_n(t_{\boldsymbol j})=
 \begin{cases}E_{\omega(\sigma)},&j_i=\sigma(i)\text{ for all }i,\\
 0,&\text{otherwise},\end{cases}\\
 K_n:T_n&\longrightarrow B_n,\qquad
 K_n(t_{\boldsymbol j})=r_n^{-1}\chi_n\prod_i d_{j_i},\\
 \varepsilon_n:T_n&\longrightarrow O,\qquad
 \varepsilon_n(t_{\boldsymbol j})=
 \begin{cases}1,&\boldsymbol j\text{ is a permutation tuple},\\
 0,&\text{otherwise}.\end{cases}
 \end{split}
\]
Writing $\tau_A$ for regular trace on $A_{\mathrm{perm}}$, we have
the equalities of actual linear maps
\begin{equation}\label{eq:compressed-trace-square}
 \tau_AS_n=\varepsilon_n=\tau_BK_n.
\end{equation}
The first follows because a coordinate idempotent has trace one;
the second is exactly the parity computation for one tuple.
The universal row-choice tensor
\[
 t(Z)=\bigotimes_{i=1}^n\Bigl(\sum_{j=1}^n Z_{ij}t_j\Bigr)
     =\sum_{\boldsymbol j}\Bigl(\prod_i Z_{ij_i}\Bigr)t_{\boldsymbol j}
 \quad\text{in }T_n\otimes_O R
\]
is sent by $S_n$ to \eqref{eq:compressed-original-u}, by $K_n$
to \eqref{eq:compressed-row-multiplier}, and by $\varepsilon_n$
to the permanent. Thus \eqref{eq:compressed-trace-square} relates
every coefficient before taking the trace of the polynomial.

The kernels and surjections in this correspondence are explicit.
The map $S_n$ is surjective, and its kernel has as basis the
nonpermutation tuples, because the remaining images are the
distinct basis idempotents. For a tuple put
\[
 p(\boldsymbol j)=\{j\in\{2,\ldots,n\}:
            j\text{ occurs an odd number of times in }\boldsymbol j\}.
\]
For each subset $U\subseteq\{2,\ldots,n\}$ choose
$\boldsymbol a(U)$ to consist first of the elements of $U$ in
increasing order and then of $n-|U|$ entries equal to one.
Its parity set is $U$. Thus
$K_n(t_{\boldsymbol a(U)})=r_n^{-1}\chi_nd_U$,
and these elements form a basis of $B_n$ because multiplication
by the unit $\chi_n/r_n$ permutes and rescales its basis.
This proves surjectivity with a specified section. The differences
\[
 t_{\boldsymbol j}-t_{\boldsymbol a(p(\boldsymbol j))}
 \quad\bigl(\boldsymbol j\ne\boldsymbol a(p(\boldsymbol j))\bigr)
\]
form a basis of $\ker K_n$: subtracting these differences reduces
an arbitrary vector to the selected representatives, whose images
are independent; each difference has coefficient one on its own
unselected basis vector and zero on the other unselected vectors,
which proves independence.

The parity proof also shows that $p(\boldsymbol j)=\{2,\ldots,n\}$
holds exactly for permutation tuples. Every other selected representative
$\boldsymbol a(U)$ is therefore nonpermutational. The regular trace
kernel in $B_n$ is the free span of $d_V$ with $V\ne\varnothing$,
by \eqref{eq:compressed-regular-trace} and the invertibility of $r_n$.
The preceding representatives consequently prove
\begin{equation}\label{eq:compressed-kernel-trace}
 K_n(\ker S_n)=\ker\tau_B.
\end{equation}

\begin{proposition}\label{prop:compressed-exact-sequence}
There is an exact sequence of $O$-modules
\[
 T_n\xrightarrow{\ (S_n,-K_n)\ }
 A_{\mathrm{perm}}\oplus B_n
 \xrightarrow{\ (a,b)\mapsto\tau_A(a)+\tau_B(b)\ }O
 \longrightarrow0.
\]
The first image has rank $n!+r_n-1$, and its kernel
$\ker S_n\cap\ker K_n$ is free of rank $n^n-n!-r_n+1$.
\end{proposition}
\begin{proof}
Equation~\eqref{eq:compressed-trace-square} says the composite is
zero. The second map is surjective because it sends
$(E_{\omega(\mathrm{id})},0)$ to one. Suppose
$\tau_A(a)+\tau_B(b)=0$. Choose $t_0\in T_n$ with $S_n(t_0)=a$.
Then $b+K_n(t_0)$ has trace
$\tau_B(b)+\varepsilon_n(t_0)=\tau_B(b)+\tau_A(a)=0$.
Equation~\eqref{eq:compressed-kernel-trace} supplies
$t_1\in\ker S_n$ such that $K_n(t_1)=-b-K_n(t_0)$.
It follows that $(S_n,-K_n)(t_0+t_1)=(a,b)$, proving exactness.

There is an explicit basis for the kernel of the second arrow:
take $(E_{\omega(\sigma)}-E_{\omega(\mathrm{id})},0)$ for
$\sigma\ne\mathrm{id}$, together with
$(0,d_V)$ for $V\ne\varnothing$ and the single vector
$(-r_nE_{\omega(\mathrm{id})},1)$. Their traces vanish.
Subtracting the nonidentity $A_{\mathrm{perm}}$ coefficients and
the nonempty $B_n$ coefficients from any trace-zero pair reduces
it to $(a_0E_{\omega(\mathrm{id})},b_0)$ with
$a_0+r_nb_0=0$, which is the last vector multiplied by $b_0$.
The same coordinates prove independence. There are $n!+r_n-1$
such vectors. Choose lifts of these basis vectors under the first
surjection just proved; their linear extension gives a section.
Then $T_n$ is the direct sum of their free span and the first
kernel. For freeness of that kernel without an unstated theorem,
decompose $T_n$ into permutation and nonpermutation basis vectors.
An element of $\ker S_n$ uses only the latter. Group those vectors
by their parity sets $U\ne\{2,\ldots,n\}$. On each group $K_n$
sends every vector to the same nonzero distinct basis multiple.
Its kernel therefore has basis the differences against the selected
representative in that group, by the argument above. The total
number is $(n^n-n!)-(r_n-1)$, as asserted.
\end{proof}

For $n\ge2$ this common-tensor construction cannot be replaced by
a linear map making either of its two surjections factor through
the other. The tuple $(1,\ldots,1)$ is nonpermutational, so its
$S_n$ image is zero whereas its $K_n$ image is the nonzero unit
$\chi_n/r_n$. This disproves a factorization $K_n=L S_n$.
For two distinct permutations, the difference of their tensor
vectors has $K_n$ image zero, since each image is $1/r_n$,
whereas its $S_n$ image is a nonzero difference of basis idempotents.
This disproves a factorization $S_n=M K_n$. The exact sequence
and coefficient trace maps above supply the proved relation despite
these specific obstructions. For $n=1$ both coefficient maps are
isomorphisms between rank-one modules, and the original full source
still has rank $3^3=27$ with its stated complement.
More explicitly, at $n=1$ the unique flagged point is
$(X_{B_{11}},X_{R_{11}},X_{C_{11}})=(1,1,1)$. Evaluation there
is a surjective $O$-algebra map $\mathcal A_3\to O=B_1$;
its kernel is the span of the other 26 point idempotents, and
it sends $u_1(Z)=Z_{11}E_{(1,1,1)}$ to $v_1(Z)=Z_{11}$.

\subsection{A multiplication obstruction and the unequal power traces}

Equality of the first traces does not give equality of their spectra.
Specialize the original coordinates to the identity matrix $I_n$.
In \eqref{eq:compressed-original-u} only $\sigma=\mathrm{id}$
survives, so $u_n(I_n)=E_{\omega(\mathrm{id})}$ is an idempotent
in the full original algebra and its permutation ideal. On the sign
side $L_i(I_n)=d_i$, and hence
\[
 v_n(I_n)=r_n^{-1}\chi_n\prod_i d_i=r_n^{-1}\chi_n^2=r_n^{-1}1.
\]
For $n\ge2$ over $\mathbb Q$ this element is not idempotent,
since $r_n^{-2}-r_n^{-1}=(1-r_n)/r_n^2\ne0$.
Every multiplicative map, whether unital or not, sends an idempotent
to an idempotent. Thus no $\mathbb Q[Z]$-linear multiplicative map
from the original full source algebra or its permutation ideal
to $B_n\otimes\mathbb Q[Z]$ sends $u_n(Z)$ to $v_n(Z)$:
such a map would specialize at $Z=I_n$ to the impossible map just
described. This is an obstruction for the stated compatible base
map. No assertion about arbitrary changes of the $Z$ coordinates
is needed. In odd characteristic the specific test may disappear
when $r_n=1$ in the field; the rational assertion retains its domain.

There is a stronger generic obstruction that covers every field $k$
of characteristic different from two, even after passing to the
rational function field $K=k(Z_{ij})$. Both source algebras and the
sign algebra split as products of copies of $K$, by their stated
point idempotents. Compose a $K$-linear multiplicative map from
either source algebra to the sign algebra with any sign-coordinate
evaluation. Each primitive source idempotent must map to an
idempotent of the field $K$, hence to zero or one: the equation
$x(x-1)=0$ in a field forces these two possibilities. Orthogonality
allows at most one of them to map to one. Therefore the image
of $u_n$ in this coordinate is either zero or one of its permutation
monomials. At a fixed sign vector $\delta$, however, the coefficient
of the nonpermutation monomial $\prod_i Z_{i1}$ in $v_n(\delta)$
is $r_n^{-1}\prod_j\delta_j\ne0$, and the coefficient of every
permutation monomial is $r_n^{-1}\ne0$. These statements follow
by selecting the indicated unique term in each row sum; all row
monomials are distinct. For $n\ge2$ this polynomial is consequently
neither zero nor a single permutation monomial. The injection
$k[Z]\hookrightarrow K$ preserves these inequalities. This proves
that no such $K$-linear multiplicative map sends $u_n$ to $v_n$,
including nonunital maps. It also proves the compatible polynomial
base-map obstruction over every such field, since a map over
$k[Z]$ would extend to $K$ by scalar extension.

For every integer $d\ge1$ the original power trace and the new one
are, respectively,
\[
 \operatorname{Tr}(m_{u_n}^d)=
       \sum_{\sigma}\prod_i Z_{i,\sigma(i)}^d,
 \qquad
 \operatorname{Tr}(m_{v_n}^d)=
       r_n^{-d}\sum_{\delta\in\Delta_n}
       \Bigl(\prod_j\delta_j\Bigr)^d
       \prod_i\Bigl(\sum_j\delta_j Z_{ij}\Bigr)^d.
\]
These follow by taking powers of the diagonal entries in their
specified point bases. At $Z=I_n$ they equal $1$ and $r_n^{1-d}$.
In particular $d=2$ distinguishes them over $\mathbb Q$ for $n\ge2$.
Their characteristic determinants at this same specialization are
$1-T$ and $(1-T/r_n)^{r_n}$, respectively: the full original
operator has one eigenvalue one and all other eigenvalues zero,
whereas the sign operator has $r_n$ copies of $1/r_n$.
The common first trace, the exact coefficient maps, and these
different higher traces are all retained simultaneously.

\subsection{The exact dimension obstruction for a fixed algebra}

Let $k$ now be any field and $D$ a finite-dimensional associative
unital $k$-algebra of dimension $q_D$. Fix $c\in D$ and elements
$a_{ij}\in D$, all independent of the polynomial coordinates $Z$.
No commutativity of $D$ is assumed. Define
\[
 A_i(Z)=\sum_{j=1}^n a_{ij}Z_{ij},\qquad
 Q_D(Z)=\operatorname{Tr}_{D\otimes k[Z]/k[Z]}
                 (m_{c A_1(Z)\cdots A_n(Z)}),
\]
where the product has the displayed order, the $Z$ coordinates are
central, and $m_b(x)=bx$ is left multiplication. Associativity
implies $m_bm_{b'}=m_{bb'}$ directly on every $x\in D$.

\begin{theorem}\label{thm:compressed-fixed-algebra-bound}
For every cut $0\le k_0\le n$, the row coefficient contraction
of $Q_D$ has rank at most $q_D$. Consequently, if
$Q_D(Z)=\operatorname{per}_n(Z)$ as polynomials, then
\[
 q_D\ge\binom n{k_0}\quad(0\le k_0\le n),\qquad
 q_D\ge\binom n{\lfloor n/2\rfloor}.
\]
\end{theorem}
\begin{proof}
Use the head and tail monomial spaces $H_{k_0},T_{k_0}$ of
Theorem~\ref{thm:per-exact-row-rank}. Define a linear map on
the dual head basis by
\[
 a:H_{k_0}^*\longrightarrow D,\qquad
 h_{\boldsymbol j}^*\longmapsto
 a_{1j_1}\cdots a_{k_0j_{k_0}},
\]
with the empty product equal to $1_D$. Define
\[
 b:D\longrightarrow T_{k_0},\qquad
 z\longmapsto\sum_{\boldsymbol l}
 \operatorname{Tr}_D\bigl(m_{c z
      a_{k_0+1,l_{k_0+1}}\cdots a_{n,l_n}}\bigr)
 t_{\boldsymbol l}.
\]
Both are $k$-linear, since multiplication by a fixed factor is
linear and the matrix trace is linear. Expanding the ordered product
in $Q_D$ shows that its coefficient at
$h_{\boldsymbol j}t_{\boldsymbol l}$ is exactly the displayed
trace with $z=a_{1j_1}\cdots a_{k_0j_{k_0}}$. Its coefficient
contraction is therefore $ba$ and has rank at most $\dim_kD=q_D$.
If it is the permanent contraction, the integral identity minor
of Theorem~\ref{thm:per-exact-row-rank} gives rank $\binom n{k_0}$
over this same field. The middle choice gives the last inequality.
\end{proof}

For any field of odd characteristic, or of characteristic zero,
\eqref{eq:compressed-row-multiplier} is a representation of this
precise form with $D=B_n\otimes k$, $c=\chi_n/r_n$,
$a_{ij}=d_j$, and $q_D=2^{n-1}$. Thus the construction gives an
upper bound $2^{n-1}$ for this algebra dimension, while the theorem
gives its universal lower bound $\binom n{\lfloor n/2\rfloor}$.
This lower bound is exponential: the binomial coefficients count
all subsets by their cardinalities and sum to $2^n$; the ratio
$\binom n{k+1}/\binom nk=(n-k)/(k+1)$ shows the middle coefficient
is maximal, so it is at least $2^n/(n+1)$.
The factorization assumes a fixed algebra with coordinate-independent
structure and coefficients. It places no restriction on representing
an arbitrary polynomial as a single unconstrained element of a
rank-one algebra; in that larger description the element itself
could already contain the permanent.

The sign algebra is smaller than the original source algebra for
every $n\ge1$, since $2^{n-1}\le3^{n-1}<3^{3n^2}$, and smaller
than its Boolean corner $2^{3n^2}$. Its stated dimension remains
exponential. Enumerating its $2^{n-1}$ signs and directly evaluating
the displayed $n$ row sums and product uses $O(n^2 2^{n-1})$
arithmetic operations: each sign requires $n^2$ signed entries,
$n(n-1)$ additions for the row sums, $n-1$ multiplications for
their product, and a bounded number of sign and scale operations;
the resulting values are summed over that many signs. These loops
give a uniform explicit algorithm. For integer $L\ge0$ and entries with
$|Z_{ij}|<2^L$, each row sum has absolute value at most $n2^L$,
and the integer signed numerator is bounded in absolute value by
$2^{n-1}n^n2^{nL}$. Its exact division by $2^{n-1}$ yields the
integer permanent by the proof above; a sufficient signed bit bound
for that numerator is
$n-1+n\lceil\log_2 n\rceil+nL+2$.
No polynomial-size representation or unrestricted circuit separation
is inferred from this compressed trace or its fixed-algebra bound.
